\section{C++}
\label{sec:c++newfeatures}
\vspace{10.5cm}
\subsection{C++ Good Practices}
\subsubsection{Rule of three/five/zero}
\nparagraph{Rule of three}
If a class requires a \textbf{user-defined destructor}, a user-defined 
\textbf{copy constructor}, or a user-defined \textbf{copy assignment 
operator}, it almost certainly requires all three.

Because C++ copies and copy-assigns objects of user-defined types in 
various situations (passing/returning by value, manipulating a 
container, etc), these special member functions will be called, 
if accessible, and if they are not user-defined, they are 
implicitly-defined by the compiler.

The implicitly-defined special member functions are typically incorrect 
if the class manages a resource whose handle is an object of non-class 
type (raw pointer, POSIX file descriptor, etc), whose destructor does 
nothing and copy constructor/assignment operator performs a 
"shallow copy" (copy the value of the handle, without duplicating the 
underlying resource).

%~ \textbf{CODE EXAMPLE HERE}

Classes that manage non-copyable resources through copyable handles 
may have to declare copy assignment and copy constructor private 
and not provide their definitions or define them as deleted. 
This is another application of the rule of three: deleting one 
and leaving the other to be implicitly-defined will most likely result 
in errors.

\nparagraph{Rule of five}
Because the presence of a user-defined destructor, copy-constructor, 
or copy-assignment operator prevents the implicit definition of the 
move constructor and the move assignment operator, \textbf{any class 
for which move semantics are desirable}, has to declare all five special 
member functions.

%~ \textbf{CODE EXAMPLE HERE}

Unlike Rule of Three, 
failing to provide move constructor and move assignment 
is usually not an error, 
but a missed optimization opportunity.

\nparagraph{Rule of zero}
Classes that have custom 
destructors, copy/move constructors or copy/move assignment operators 
should deal exclusively with ownership 
(which follows from the Single Responsibility Principle). 
Other classes should not have custom 
destructors, copy/move constructors or copy/move assignment operators.
(see~\cite{fernandesruleofzero})

This rule also appears 
in the C++ Core Guidelines as C.20: 
If you can avoid defining default operations, do.

When a \textbf{base class is intended for polymorphic use}, 
its destructor may have to be declared public and virtual. 
This blocks implicit moves (and deprecates implicit copies), 
and so the special member functions have to be declared as defaulted
(see~\cite{meyersconcernaboutruleof0})

However, this makes the class prone to slicing, 
which is why polymorphic classes often 
define copy as deleted 
(see~\cite{cpp-core-guidelines-C67}), 
which leads to the following generic wording for the Rule of Five:
If you define or =delete any default operation, 
define or =delete them all
(see~\cite{cpp-core-guidelines-C21})

\subsubsection{Const Correctness}
Const Correctness consists in using the keyword \texttt{const} to 
prevent \texttt{const} objects from being mutated.

\nparagraph{const and pointer combinations}
Reading it right-to-left
\begin{table}[h!t]
    \begin{tabular}{l|r}
        \texttt{const X* p} & p is a pointer to an X that is constant.\\
        \hline
        \texttt{X const* p} & equivalent to \texttt{const X* p}.\\
        \hline
        \texttt{X* const p} & p is a \texttt{const} pointer to an X that 
                              is not \texttt{const}.\\
        \hline
        \texttt{const X* const p} &  p is a \texttt{const} pointer to an 
                                     X that is \texttt{const}.\\
    \end{tabular}
\end{table}

The second combination is equivalent to the first one because of the 
\textbf{\texttt{const}-on-the-right style}, which require to write the 
\texttt{const} specifier on the right of the object it makes 
\texttt{const}. This style is also called "consistent \texttt{const}" 
since it is consistent with the way of defining a const member function. 
Moreover, it is also more consistent with pointer declarations.
\nparagraph{const and reference combinations}
Reading it right-to-left
\begin{table}[h!t]
    \begin{tabular}{l|r}
        \texttt{const X\& r} & r aliases an X object, which can not be changed via r \\
        \hline
        \texttt{X const\& r} & equivalent to \texttt{const X\& r} \\
        \hline
        \texttt{X\& const r} & nonsense functionally equivalent to \texttt{X\& r}\\
        \hline
        \texttt{const X\& const r} & nonsense functionally equivalent to \texttt{const X\& r}\\       
    \end{tabular}
\end{table}
The equivalence between the second and the first combination is due 
to the same reason valid for the pointers.
\newline
The last two combinations are nonsense and should be avoided. They are 
redundant as a reference is always constant, in the sense that you can 
never reseat a reference to make it refer to a different object.

\nparagraph{Logical and Physical state}
When reasoning in terms of const-correctness, it's important to 
distinguish the logical and the physical state of an object.
\newline
On the inside, an object has a \textbf{physical state}, also called 
\textbf{concrete} or \textbf{bitwise} state. This state is the one 
visible to the programmer who implements the software. From the outside, 
the client has an access to the object restricted to public member 
functions and friend functions. What is visible is the \textbf{logical} 
or \textbf{abstract} or \textbf{meaningwise state}.
\newline
A lookup functionality for a collection-object is a good example to 
show how there could be bits in the object's physical state that haev 
no corresponding elements in the object's logical state. The lookup 
functionality might improve the performance using a cache to store 
the already looked-up items. The implementation of the cache is part of 
the object's physical state but the implementation is internal and 
the details will probably not be exposed to the users, i.e. it is not 
part of the object's logical state.

\nparagraph{\texttt{const}ness of \texttt{public} member functions}
The constness of a public method must make sense to the object's users, 
and those users can see only the object's logical state. Therefore, in 
this context the \texttt{const}ness to take care of is the logical one. 
\begin{itemize}    
    \item If a method changes any part of the object’s logical state, 
          it logically is a mutator; it should not be const even if 
          (as actually happens!) the method doesn’t change any physical 
          bits of the object’s concrete state.
    \item Conversely, a method is logically an inspector and should be 
          const if it never changes any part of the object’s logical 
          state, even if (as actually happens!) the method changes 
          physical bits of the object’s concrete state.
\end{itemize}

\nparagraph{return-by-reference and \texttt{const} member function}
A \texttt{const} member function must not change nor allow its caller to 
change the logical state of the \texttt{this} object it's been called 
on.
\newline

\nparagraph{\texttt{const} overloading}
\nparagraph{\texttt{mutable} keyword}
\nparagraph{Error converting \texttt{F**} into \texttt{const F**}}

\subsection{C++ General Knowledge}
\subsubsection{Classes and Inheritance}
This is section is based upon the \href{https://isocpp.org/wiki/faq/basics-of-inheritance}{isocpp FAQ about inheritance}
\nparagraph{Virtual member function}
A \href{https://en.cppreference.com/w/cpp/language/virtual}{virtual member function} 
is the way of C++ to allow derived classes to replace the implementation 
provided by the base class. The compiler makes sure the correct 
derived class' implementation is called, even when the object is 
accessed through a pointer to the base class.
\newline
The reason why the member function are not \texttt{virtual} by default 
is that many classes are not designed to be used as a base class. 
Moreover, the \texttt{virtual} function call mechanism needs requires 
an overhead of memory, typically one word per object.
\nparagraph{Dynamic Binding and Static Typing}
\textbf{Static typing} means that the legality of a member function 
invocation is checked as early as possible, i.e. at compile time by the 
compiler itself. In this kind of check, the compiler uses the 
\textbf{static type} of the pointer to determine whether the member 
function invocation is correct, meaning that the type pointed by the 
pointer can handle the member function.
\textbf{Dynamic binding} is the \texttt{virtual} function call 
mechanism, which checks the correctness of the function call at the last 
possible moment, i.e. at run time and based on the dynamic type of 
the object pointed to.

\subsubsection{Member Function call: \texttt{virtual} vs. non-\texttt{virtual}}
The practical difference is that non-\texttt{virtual} member functinos 
are resolved statically, i.e. the correct member function is selected 
statically \textbf{depending on the type of the pointer or reference to 
the object}.
\newline
On the other side, a \texttt{virtual} member function call is resolved 
dynamically 
\subsubsection{Virtual Member Function Call - Overhead Estimation}
\subsubsection{Calling a Base Class \texttt{virtual} member function}
\subsubsection{Heterogeneous list of objects}
\subsubsection{\texttt{virtual} destructor}
\subsubsection{\texttt{virtual} constructor}


\subsection{C++11 New Features}
\subsubsection{Value Categories}
\href{https://en.cppreference.com/w/cpp/language/value_category}{cppreference documentation} 
about value categories.
\newline    
Due to the expressiveness of the C++ language, the terms \textbf{lvalue} 
and \textbf{rvalue} references evolved in a more complex set.
\nparagraph{glvalues}
\textbf{G}eneralized lvalue is used to refer to something that could be 
either a \textbf{lvalue} or \textbf{xvalue}.
\nparagraph{rvalues}
This term is inherited from C, in which \textbf{rvalues} can be on the 
\textbf{r}ight side of an assignment. In C++ this term can refer to 
things that are either \textbf{xvalues} and \textbf{prvalues}.
\nparagraph{lvalues}
This term is also inherited from C, in which \textbf{lvalues} can be 
on the \textbf{l}eft side of an assignment. Therefore, the simplest one 
is the name of a variable in a declaration.
\newline
Any expression resulting in an \textbf{lvalue reference} is also an 
\textbf{lvalue}.
\begin{lstlisting}[language=C+++]
int v;          // v is an lvalue of type int
A*  p1;         // *p1 is an lvalue
A&  r1 = v1;    // r1 is an lvalue
A&  f();        // f() is an lvalue
\end{lstlisting}

Accessing memebers of objects through pointers or references yields 
\textbf{lvalues}. For example, with the following code
\begin{lstlisting}[language=C+++]
struct Foo
{
    Bar a;
    Bar b;
};

Foo& f2();
Foo* p2;
Foo  v2;
\end{lstlisting}
\texttt{f2().a}, \texttt{p2->b} and \texttt{v2.a} are all 
\textbf{lvalues} if type A.
\newline
\textbf{String literals} are \textbf{lvalues} of type array.
\newline
To conclude, a \textbf{named object} declared with an 
\textbf{rvalue reference} is an \textbf{lvalue}. The name rvalue 
reference indicates that it can bind to an \textbf{rvalue}, and once 
you've declared a variable and given it a name it is an \textbf{lvalue}.
\begin{lstlisting}[language=C+++]
void Foo(A&& bar)
{
    // within Foo, bar is an lvalue 
    // which will only bind to rvalues
}
\end{lstlisting}

\nparagraph{xvalues}
This term indicates e\textbf{X}piring value: an unnamed object that 
will be destroyed soon. \textbf{xvalues} can be treated either as 
\textbf{glvalues} or \textbf{rvalues}, based on the context.
\newline
Usually, \textbf{xvalues} only arise through \textbf{explicit casts} and 
\textbf{function calls}. If an expression is cast to an rvalue reference 
to some type T then the result is an xvalue of type T 
\begin{lstlisting}[language=C+++]
static_cast<A&&>(v1) // yields an xvalue of type A
\end{lstlisting}
If the return type of a function is an \textbf{rvalue reference} to some 
type \textbf{T}, the result is an \textbf{xvalue} of type \textbf{T}.
For example, the \texttt{std::move} is declared as
\begin{lstlisting}[language=C+++]
template <typename T>
constexpr remove_reference_t<T>&& move(T&& t) noexcept;
\end{lstlisting}
\texttt{std::move(v1)} is an \textbf{xvalue} of type \texttt{A}. 
The type deduction rules deduce \textbf{T} to be \texttt{A\&} 
because \texttt{v1} is an \textbf{lvalue}.

\nparagraph{prvalues}
\textbf{P}ure \textbf{rvalue} is an \textbf{rvalue} which is not an  
\textbf{xvalue}
\newline
Literals other than string literals are \textbf{pvalues}. For example, 
$42$ is a \textit{prvalue} of type \texttt{int} and $3.141f$ is a 
\textit{prvalue} of type \texttt{float}
\newline
\textbf{Temporaries} are \textbf{prvalues}. For example any temporaries 
created as a result of an implicit conversion is also a \textbf{prvalue}. 
In the following code
\begin{lstlisting}[language=C+++]
int consume_string(std::string&& s);
int i = consume_string("hello");
\end{lstlisting}
the string literal in the function call will implicitly convert to a 
temporary of type 
\newline 
\texttt{std::string}, which can bind to the rvalue 
reference used for the function parameter since the temporary is a 
\textit{prvalue}
\nparagraph{Reference Binding}
\href{https://en.cppreference.com/w/cpp/language/reference}{cppreference documentation} 
about reference declaration.
\newline
The main difference between \textbf{lavalues} and \textbf{rvalues} is 
the way they bind to references. The differences in type deduction rules 
can have a big impact too.
\nsubparagraph{lvalue references}
These references are declared with a single ampersand $\&$.
\newline
A \textbf{non-const lvalue reference will only bind to non-const 
lvalues} of the same type, or a class derived from the referenced type
\begin{lstlisting}[language=C+++]
class C : public A{};

int i = 42;
A a;
B b;
C c;
const A ca{};

A&   r1  = a;
A&   r2  = c;
A&   r3  = b;     // error, wrong type
int& r4  = i;
int& r5  = 42;    // error, cannot bind rvalue
A&   r6  = ca;    // error, cannot bind const object
                  // to non const ref
A&   r7  = r1;
A&   r8  = A();   // error, cannot bind rvalue
A&   r9  = B().a; // error, cannot bind rvalue
A&   r10 = C();   // error, cannot bind rvalue
\end{lstlisting}
A \textbf{const lvalue reference can also bind to rvalues}. The 
object bound to the reference must have the same type as the referenced 
type or a class derived from the referenced type. It is possible to 
bind both const and non-const values to a const lvalue reference.
\begin{lstlisting}[language=C+++]
const A&   cr1  = a;
const A&   cr2  = c;
const A&   cr3  = b;     // error, wrong type
const int& cr4  = i;
const int& cr5  = 42;    // rvalue can bind OK
const A&   cr6  = ca;    // OK, can bind const object to const ref
const A&   cr7  = cr1;
const A&   cr8  = A();   // OK, can bind rvalue
const A&   cr9  = B().a; // OK, can bind rvalue
const A&   cr10 = C();   // OK, can bind rvalue
const A&   cr11 = r1;
\end{lstlisting}
Binding a temporary object (i.e. a \textbf{prvalue}) to a const lvalue 
reference \textbf{extends the lifetime} of that temporary to the 
lifetime of the reference. In the previous example, it is possible 
to use \texttt{r8}, \texttt{r9} and \texttt{r10} later in the code.
\newline
Lifetime extension does not extend to references initialized from the 
first reference. If a function parameter is a const lvalue reference 
and gets bound to a temporary passed to the function call, the temporary 
is destroyed when the function returns. This happens also if the 
reference was stored in a longer-lived variable, such as a member of a 
newly constructed object.
\newline
\textbf{volatile} and \textbf{const volatile} lvalue references are much 
less used but behave as expected: volatile T\& bind to a volatile or 
non-volatile, non-const lvalue of type T or a class derived from T. 
However, volatile const lvalue references do not bind to rvalues.

\nsubparagraph{rvalue references}
An \textbf{rvalue reference} can be used to extend the lifetime of 
temporary objects, making the referenced object modifiable through 
them. The reference must be const in order to bind to a const object, 
though const rvalue reference are much rarer.
\begin{lstlisting}[language=C+++]
const A make_const_A();
A&&       rr1  = a;                // error, cannot bind lvalue 
                                   // to rvalue reference
A&&       rr2  = A();
A&&       rr3  = B();              // error, wrong type
int&&     rr4  = i;                // error, cannot bind lvalue
int&&     rr5  = 42;
A&&       rr6  = make_const_A();   // error, cannot bind 
                                   //        const object 
                                   //        to non const ref
const A&& rr7  = A();
const A&& rr8  = make_const_A();
A&&       rr9  = B().a;
A&&       rr10 = C();
A&&       rr11 = std::move(a);     // std::move returns an rvalue
A&&       rr12 = rr11;             // error rvalue references 
                                   //       are lvalues
\end{lstlisting}
rvalue references extend the lifetime of temporary objects in the same 
way const lvalue references do, therefore temporaries associated with 
\texttt{rr2}, \texttt{rr5}, \texttt{rr7}, \texttt{rr8}, \texttt{rr9} 
and \texttt{rr10} remain valid until the corresponding references are 
destroyed. 

\nsubparagraph{Reference Collapsing}
In C++ it is possible to form reference to reference in 
\textbf{templates} and \texttt{typedef} expressions. In this context, 
the refefrence collapsing rule apply: rvalue reference to rvalue 
reference collapses to an rvalue reference, all the others collapse to 
lvaue references.
\begin{lstlisting}[language=C+++]
typedef int&  lref;
typedef int&& rref;
int n;
 
lref&  r1 = n; // type of r1 is int&
lref&& r2 = n; // type of r2 is int&
rref&  r3 = n; // type of r3 is int&
rref&& r4 = 1; // type of r4 is int&&
\end{lstlisting}

\nsubparagraph{Implicit conversion}
A reference-to-A can be initialized with any D object which has an 
inplicit conversion operator that returns \texttt{A\&}. In this case 
the reference is bound to the result of the conversion operator.
\newline
A const lvalue-reference-to-E will bind only to objects implicitly 
convertible to E. The reference is bound to the temporary resulting 
from the conversion.

\nsubparagraph{General properties}
In general, \textbf{rvalues} can't be modified and their address can't 
be taken.
\newline
However, as long as class \texttt{X} has a member function 
named \texttt{do\char`_something}, the expression 
\texttt{X().do\char`_something} is valid.
\newline
\textbf{ref qualifiers} along with \texttt{const} can be used to tag 
different overloads of class member functions, to indicate they can be 
applied to only \textit{lvalues} or to only \textit{rvaules}.
\newline
When the type of a variable or function template parameter is deduce 
from its initializer, whether the initializer is an lvalue or rvalue 
can affect the deduced type.

\subsubsection{Move Semantic}
Move semantics is a new way of moving resources around in an optimal way 
by \textbf{avoiding unnecessary copies} of temporary objects. It is 
based on rvalue references and plays an \textbf{important role in 
exception safety} for C++. 

%~ \subsubsection{Universal Reference}
%~ \href{https://en.cppreference.com/w/cpp/language/reference#Forwarding_references}{cppreference documentation}
%~ about universal reference
%~ \newline
\subsubsection{Perfect Forwarding}
\href{https://en.cppreference.com/w/cpp/utility/forward}{cppreference documentation} 
about \textbf{forward} utility.
\newline
\href{https://en.cppreference.com/w/cpp/language/reference#Forwarding_references}{cppreference documentation} 
about \textbf{forwarding references}.
\newline
Perfect forwarding is an important C++0x feature built on top of rvalue 
references. It provides the programmer with the capability automatically 
apply the move semantics, even when there are intervening function calls 
to separate the soure and the destination of a move.
\newline
Examples include constructors and setter functions, that can forward 
arguments they receive directly to the data member they are initializing 
or setting. Standard Library functions like \texttt{make\char`_shared} 
"perfect-forwards" its arguments to the class constructor of whatever 
object the to-be-created \texttt{shared\char`_ptr} is to point to.
\newline
The idiomatic expression of the perfect forwarding make use of the 
\texttt{std::forward} 
\begin{lstlisting}[language=C+++]
template <typename T>
void wrapper(T&& arg) 
{
    wrapped(std::forward<T1>(t1), std::forward<T2>(t2));
}
\end{lstlisting}

\subsubsection{auto}
\subsubsection{decltype}
\subsubsection{const-expr}

\subsection{C++23 New Features}
\subsubsection{Containers}
With C++23, we have 12 associative Containers:
\begin{table}[h!t]
    \begin{tabular}{|l|l|}
    \textbf{C++98} & \texttt{std::map}, \texttt{std::set}, \texttt{std::multimap}, \texttt{std::multiset} \\
    \hline
    \multirow{2}{*}{\textbf{C++11}} 
    & \texttt{std::unordered\char`_map}, \texttt{std::unordered\char`_set} \\
    & \texttt{std::unordered\char`_multimap}, \texttt{std::unordered\char`_multiset} \\
    \hline
    \multirow{2}{*}{\textbf{C++23}} 
    & \texttt{std::flat\char`_map}, \texttt{std::flat\char`_set},\\
    & \texttt{std::flat\char`_multimap}, \texttt{std::flat\char`_multiset} \\
    \end{tabular}
\end{table}
\newpage
\begin{table}[h!t]
    \begin{tabular}{l|l|l|l|l}
        \multirow{2}{*}{Associative Container} & \multirow{2}{*}{Sorted} & \multirow{2}{*}{Value} & Several Identical & \multirow{2}{*}{Access Time} \\
                                               &        &       & keys possible & \\
        \hline
        \texttt{std::set} & yes & no & no & logarithmic \\
        \texttt{std::unordered\char`_set} & no & no & no & constant \\
        \texttt{std::map} & yes & yes & no & logarithmic \\
        \texttt{std::unordered\char`_map} & no & yes & no & constant \\
        \texttt{std::multiset} & yes & no & yes & logarithmic\\
        \texttt{std::unordered\char`_multiset} & no & no & yes & constant\\
        \texttt{std::multimap} & yes & yes & yes & logarithmic\\
        \texttt{std::unordered\char`_multimap} & no & yes & yes & constant\\
    \end{tabular}
\end{table}
\newpage{}

